\documentclass{article}
\usepackage{amsmath,amssymb,amsthm}
\usepackage{algorithm}
\usepackage{algorithmic}
\usepackage{enumitem}
\usepackage{hyperref}
\newtheorem{theorem}{Theorem}
\newtheorem{lemma}{Lemma}
\newtheorem{assumption}{Assumption}
\newcommand{\EE}{\mathbb{E}}
\newcommand{\RR}{\mathbb{R}}
\newcommand{\norm}[1]{\left\lVert#1\right\rVert}
\newcommand{\inner}[2]{\langle #1,#2\rangle}
\begin{document}

\section*{Algorithm Name}
\textbf{Federated Averaging (FedAvg) with Local SGD and Client Drift $\delta$}

\section*{Mathematical Formulation}
Let $f_k:\RR^d\rightarrow\RR$ be the local loss on client $k\in\{1,\dots ,K\}$ and 
\[
f(w)=\frac1K\sum_{k=1}^{K}f_k(w)
\]
the global objective.  
Denote by $w_t\in\RR^d$ the global model after the $t$‑th communication round.  
Each round consists of $E$ local SGD steps on every participating client.

\begin{itemize}
\item \textbf{Local SGD (client $k$, epoch $e$).}
\[
\begin{aligned}
w_{t}^{k,0} &\gets w_t,\\
g_{t}^{k,e} &\gets \nabla f_k\!\bigl(w_{t}^{k,e-1};\xi_{t}^{k,e}\bigr),\\
w_{t}^{k,e} &\gets w_{t}^{k,e-1}-\eta\, g_{t}^{k,e},
\qquad e=1,\dots ,E,
\end{aligned}
\]
where $\xi_{t}^{k,e}$ is a random mini‑batch drawn on client $k$ and $\eta>0$ is the learning rate.

\item \textbf{Server aggregation.}
After $E$ local steps each client sends $w_{t}^{k,E}$ to the server and the server forms
\[
w_{t+1}\;=\;\frac1K\sum_{k=1}^{K} w_{t}^{k,E}.
\]

\item \textbf{Client drift.}
Define the drift vector
\[
\delta_k(w)\;:=\;\nabla f_k(w)-\nabla f(w),
\qquad 
\Delta_t\;:=\;\frac1K\sum_{k=1}^{K}\delta_k\bigl(w_{t}^{k,E}\bigr).
\]
The term $\Delta_t$ captures heterogeneity among clients.
\end{itemize}

\section*{Pseudocode}
\begin{algorithm}[H]
\caption{FedAvg with $K$ clients, $E$ local epochs}
\label{alg:fedavg}
\begin{algorithmic}[1]
\STATE \textbf{Input:} initial model $w_0$, learning rate $\eta$, local epochs $E$, total rounds $R$.
\FOR{$t=0,\dots ,R-1$}
    \FOR{each client $k=1,\dots ,K$ \textbf{in parallel}}
        \STATE $w_{t}^{k,0}\gets w_t$
        \FOR{$e=1$ \TO $E$}
            \STATE Sample mini‑batch $\xi_{t}^{k,e}$
            \STATE $g_{t}^{k,e}\gets \nabla f_k\!\bigl(w_{t}^{k,e-1};\xi_{t}^{k,e}\bigr)$
            \STATE $w_{t}^{k,e}\gets w_{t}^{k,e-1}-\eta\, g_{t}^{k,e}$
        \ENDFOR
        \STATE Send $w_{t}^{k,E}$ to server
    \ENDFOR
    \STATE $w_{t+1}\gets \frac1K\sum_{k=1}^{K} w_{t}^{k,E}$
\ENDFOR
\STATE \textbf{Output:} $w_R$
\end{algorithmic}
\end{algorithm}

\section*{Dry Run Example}
Consider the scalar quadratic $f(w)=(w-3)^2$, so $\nabla f(w)=2(w-3)$.  
Take $K=1$ (single client), $E=1$, learning rate $\eta=0.1$, and $w_0=0$.

\begin{enumerate}
\item \textbf{Round $t=0$ (local step $e=1$).}
\[
g_0^{1,1}= \nabla f( w_0 ) = 2(0-3) = -6 .
\]
\[
w_{0}^{1,1}= w_0 - \eta g_0^{1,1}=0-0.1(-6)=0.6 .
\]
\[
f(w_{0}^{1,1}) = (0.6-3)^2 = 5.76 .
\]

\item \textbf{Aggregation (only one client).}
\[
w_1 = w_{0}^{1,1}=0.6 .
\]

\item \textbf{Round $t=1$.}
\[
g_1^{1,1}=2(0.6-3) = -4.8,
\qquad
w_{1}^{1,1}=0.6-0.1(-4.8)=1.08,
\]
\[
f(w_{1}^{1,1})=(1.08-3)^2=3.6864 .
\]

\item \textbf{Round $t=2$.}
\[
g_2^{1,1}=2(1.08-3)=-3.84,
\qquad
w_{2}^{1,1}=1.08-0.1(-3.84)=1.464,
\]
\[
f(w_{2}^{1,1})=(1.464-3)^2=2.3689 .
\]
\end{enumerate}
The objective value decreases at each round, illustrating the algorithm’s behaviour.

\newpage
\section*{Assumptions}
\begin{assumption}[Smoothness]
Each local loss $f_k$ is $L$‑smooth:
\[
\forall w,w'\in\RR^d,\qquad 
f_k(w')\le f_k(w)+\inner{\nabla f_k(w)}{w'-w}
+\frac{L}{2}\norm{w'-w}^2 .
\tag{A1}
\]
\end{assumption}

\begin{assumption}[Bounded stochastic variance]
For any $w$ and any client $k$,
\[
\EE_{\xi\sim\mathcal{D}_k}\!\bigl[\norm{\nabla f_k(w;\xi)-\nabla f_k(w)}^2\bigr]\le\sigma^2 .
\tag{A2}
\]

\end{assumption}

\begin{assumption}[Client heterogeneity (drift bound)]
Define $\delta_k(w)=\nabla f_k(w)-\nabla f(w)$. There exists $\zeta\ge0$ such that
\[
\frac1K\sum_{k=1}^{K}\norm{\delta_k(w)}^2\le \zeta^2 ,\qquad\forall w\in\RR^d .
\tag{A3}
\]
\end{assumption}

\begin{assumption}[Learning rate]
The stepsize satisfies
\[
0<\eta\le\frac{1}{L\sqrt{E}} .
\tag{A4}
\]
\end{assumption}

\section*{Main Convergence Theorem}
\begin{theorem}\label{thm:main}
Assume (A1)–(A4) and let $T=RE$ be the total number of stochastic gradient evaluations per client.  
Define the random iterate
\[
\widetilde w \;\text{chosen uniformly from }\{w_t\}_{t=0}^{R-1}.
\]
Then
\[
\EE\!\bigl[\norm{\nabla f(\widetilde w)}^2\bigr]
\;\le\;
\underbrace{\frac{2\bigl(f(w_0)-f^\star\bigr)}{\eta R E}}_{\text{optimization term}}
\;+\;
\underbrace{L\eta\bigl(\sigma^2+E\zeta^2\bigr)}_{\text{variance \& drift term}} .
\tag{T1}
\]
Choosing $\eta = \Theta\!\bigl(\tfrac{1}{\sqrt{R}}\bigr)$ yields
\[
\EE\!\bigl[\norm{\nabla f(\widetilde w)}^2\bigr]=O\!\Bigl(\frac{1}{\sqrt{R}}\Bigr)
=O\!\Bigl(\frac{1}{\sqrt{T/E}}\Bigr) .
\]
\end{theorem}

\section*{Theoretical Properties}
\begin{itemize}
\item \textbf{Stability.} The bound (T1) holds for any $E\ge1$ as long as (A4) is satisfied; larger $E$ increases the drift contribution $L\eta E\zeta^2$.
\item \textbf{Hyper‑parameter sensitivity.}
    \begin{enumerate}[label=(\alph*)]
    \item \emph{Stepsize $\eta$.} Too large $\eta$ violates (A4) and the smoothness inequality, breaking the descent guarantee.
    \item \emph{Local epochs $E$.} Increasing $E$ reduces communication but linearly inflates the drift term; optimal $E$ balances $1/(RE)$ and $E$ in (T1).
    \end{enumerate}
\item \textbf{Comparison to baseline SGD.}
    Setting $E=1$ recovers the classic non‑convex SGD rate $O(1/\sqrt{RT})$.
\end{itemize}

\section*{Preliminaries (Definitions and Lemmas)}
\begin{lemma}[Smoothness inequality]\label{lem:smooth}
If $f$ is $L$‑smooth then for any $w,w'$,
\[
f(w')\le f(w)+\inner{\nabla f(w)}{w'-w}+\frac{L}{2}\norm{w'-w}^2 .
\]
\end{lemma}
\begin{proof}
Standard from $L$‑smoothness; see (A1). \qedhere
\end{proof}

\begin{lemma}[Unbiased stochastic gradient]\label{lem:unbiased}
For any $k$ and $w$,
\[
\EE_{\xi}[g_{t}^{k,e}\mid w_{t}^{k,e-1}]=\nabla f_k\!\bigl(w_{t}^{k,e-1}\bigr).
\]
\end{lemma}
\begin{proof}
Direct from definition of $g_{t}^{k,e}$ and (A2). \qedhere
\end{proof}

\section*{Main Proof (Step‑by‑Step Derivation)}
We prove (T1) by bounding the decrease of the global objective after one communication round.
All expectations are taken w.r.t. the randomness of the stochastic gradients and the uniform choice of $\widetilde w$.

\subsection*{Step 1: Smoothness of the global loss}
From Lemma~\ref{lem:smooth} applied to $f$ and $w_{t+1}$,
\[
\begin{aligned}
f(w_{t+1})
&\le f(w_t)+\inner{\nabla f(w_t)}{w_{t+1}-w_t}
+\frac{L}{2}\norm{w_{t+1}-w_t}^2 .
\tag{1}
\end{aligned}
\]

\subsection*{Step 2: Express $w_{t+1}-w_t$ via local updates}
Recall $w_{t+1}=K^{-1}\sum_{k}w_{t}^{k,E}$ and $w_{t}^{k,0}=w_t$. Unrolling the $E$ SGD steps,
\[
w_{t}^{k,E}=w_t-\eta\sum_{e=1}^{E} g_{t}^{k,e}.
\]
Hence
\[
w_{t+1}-w_t
=-\frac{\eta}{K}\sum_{k=1}^{K}\sum_{e=1}^{E} g_{t}^{k,e}.
\tag{2}
\]

\subsection*{Step 3: Plug (2) into (1) and take expectation}
Using (2) in (1) gives
\[
\begin{aligned}
f(w_{t+1})
&\le f(w_t)
-\eta\Big\langle \nabla f(w_t),\;
\frac1K\sum_{k=1}^{K}\sum_{e=1}^{E} g_{t}^{k,e}\Big\rangle
+\frac{L\eta^{2}}{2}\norm{\frac1K\sum_{k=1}^{K}\sum_{e=1}^{E} g_{t}^{k,e}}^{2}.
\tag{3}
\end{aligned}
\]

Take full expectation $\EE[\cdot]$ conditioned on $w_t$.
Using Lemma~\ref{lem:unbiased} and the definition of drift $\delta_k$,
\[
\EE\!\bigl[g_{t}^{k,e}\mid w_t\bigr]
= \nabla f_k\!\bigl(w_t\bigr)
= \nabla f(w_t)+\delta_k(w_t).
\tag{4}
\]

Therefore,
\[
\EE\!\Bigl[\frac1K\sum_{k}\sum_{e} g_{t}^{k,e}\Bigm| w_t\Bigr]
=E\,\nabla f(w_t)+E\,\underbrace{\frac1K\sum_{k}\delta_k(w_t)}_{\Delta_t}.
\tag{5}
\]

\subsection*{Step 4: Bound the inner product term}
From (3) and (5),
\[
\begin{aligned}
\EE\!\bigl[f(w_{t+1})\mid w_t\bigr]
&\le f(w_t)
-\eta E\;\norm{\nabla f(w_t)}^{2}
-\eta E\;\inner{\nabla f(w_t)}{\Delta_t}
+\frac{L\eta^{2}}{2}\EE\!\Bigl[\norm{S_{t}}^{2}\Bigr],
\end{aligned}
\tag{6}
\]
where we introduced the shorthand
\[
S_{t}:=\frac1K\sum_{k=1}^{K}\sum_{e=1}^{E} g_{t}^{k,e}.
\]

\subsection*{Step 5: Control the drift inner product}
Using Cauchy–Schwarz and (A3):
\[
\bigl|\inner{\nabla f(w_t)}{\Delta_t}\bigr|
\le \norm{\nabla f(w_t)}\;\norm{\Delta_t}
\le \frac12\Bigl(\norm{\nabla f(w_t)}^{2}+\norm{\Delta_t}^{2}\Bigr).
\tag{7}
\]
Taking expectation and applying (A3) yields
\[
\EE\!\bigl[\norm{\Delta_t}^{2}\mid w_t\bigr]\le \zeta^{2}.
\tag{8}
\]

Insert (7) into (6):
\[
\EE\!\bigl[f(w_{t+1})\mid w_t\bigr]
\le f(w_t)
-\frac{\eta E}{2}\norm{\nabla f(w_t)}^{2}
+\frac{\eta E}{2}\zeta^{2}
+\frac{L\eta^{2}}{2}\EE\!\bigl[\norm{S_{t}}^{2}\mid w_t\bigr].
\tag{9}
\]

\subsection*{Step 6: Bounding $\EE[\norm{S_{t}}^{2}]$}
Write $g_{t}^{k,e}= \nabla f_k(w_t)+\underbrace{(g_{t}^{k,e}-\nabla f_k(w_t))}_{\varepsilon_{t}^{k,e}}$.
Then
\[
S_t = \frac{E}{K}\sum_{k=1}^{K}\nabla f_k(w_t)+\frac1K\sum_{k=1}^{K}\sum_{e=1}^{E}\varepsilon_{t}^{k,e}.
\]
Hence
\[
\begin{aligned}
\EE\!\bigl[\norm{S_t}^{2}\mid w_t\bigr]
&= \EE\!\Bigl[\Bigl\| \frac{E}{K}\sum_{k}\nabla f_k(w_t)
          +\frac1K\sum_{k,e}\varepsilon_{t}^{k,e}\Bigr\|^{2}\Bigr] \\
&\le 2\underbrace{\norm{\frac{E}{K}\sum_{k}\nabla f_k(w_t)}^{2}}_{\text{(i)}}
   +2\underbrace{\EE\!\Bigl[\norm{\frac1K\sum_{k,e}\varepsilon_{t}^{k,e}}^{2}\Bigr]}_{\text{(ii)}} .
\end{aligned}
\tag{10}
\]

\paragraph{Term (i).}
Because $\frac1K\sum_{k}\nabla f_k(w_t)=\nabla f(w_t)$,
\[
\text{(i)} = E^{2}\norm{\nabla f(w_t)}^{2}.
\tag{11}
\]

\paragraph{Term (ii).}
The noise terms are zero‑mean and independent across $(k,e)$, and each satisfies (A2):
\[
\EE\!\bigl[\norm{\varepsilon_{t}^{k,e}}^{2}\mid w_t\bigr]\le\sigma^{2}.
\]
Therefore,
\[
\text{(ii)}\le \frac{1}{K^{2}}\sum_{k=1}^{K}\sum_{e=1}^{E}\EE\!\bigl[\norm{\varepsilon_{t}^{k,e}}^{2}\bigr]
= \frac{E\sigma^{2}}{K}.
\tag{12}
\]

Plug (11)–(12) into (10):
\[
\EE\!\bigl[\norm{S_t}^{2}\mid w_t\bigr]
\le 2E^{2}\norm{\nabla f(w_t)}^{2}+ \frac{2E\sigma^{2}}{K}.
\tag{13}
\]

\subsection*{Step 7: Substitute (13) into (9)}
\[
\begin{aligned}
\EE\!\bigl[f(w_{t+1})\mid w_t\bigr]
&\le f(w_t)
-\frac{\eta E}{2}\norm{\nabla f(w_t)}^{2}
+\frac{\eta E}{2}\zeta^{2}\\
&\qquad
+\frac{L\eta^{2}}{2}\Bigl(2E^{2}\norm{\nabla f(w_t)}^{2}
      +\frac{2E\sigma^{2}}{K}\Bigr)\\[1mm]
&= f(w_t)
-\underbrace{\Bigl(\frac{\eta E}{2}-L\eta^{2}E^{2}\Bigr)}_{\alpha}\norm{\nabla f(w_t)}^{2}
+\frac{\eta E}{2}\zeta^{2}
+ \frac{L\eta^{2}E\sigma^{2}}{K}.
\end{aligned}
\tag{14}
\]

\subsection*{Step 8: Ensure a positive coefficient $\alpha$}
Using the stepsize restriction (A4), $\eta\le\frac{1}{L\sqrt{E}}$, we have
\[
L\eta^{2}E^{2}\le \eta E/2,
\]
so $\alpha\ge \frac{\eta E}{4}>0$.
Thus,
\[
\EE\!\bigl[f(w_{t+1})\mid w_t\bigr]
\le f(w_t)
-\frac{\eta E}{4}\norm{\nabla f(w_t)}^{2}
+\frac{\eta E}{2}\zeta^{2}
+ \frac{L\eta^{2}E\sigma^{2}}{K}.
\tag{15}
\]

\subsection*{Step 9: Telescope over $t=0,\dots ,R-1$}
Take total expectation and sum (15):
\[
\begin{aligned}
\EE[f(w_R)] - f(w_0)
&\le
-\frac{\eta E}{4}\sum_{t=0}^{R-1}\EE\!\bigl[\norm{\nabla f(w_t)}^{2}\bigr]
+ \frac{\eta E R}{2}\zeta^{2}
+ \frac{L\eta^{2}E R\sigma^{2}}{K}.
\end{aligned}
\tag{16}
\]

Rearrange to isolate the average squared gradient:
\[
\frac{1}{R}\sum_{t=0}^{R-1}\EE\!\bigl[\norm{\nabla f(w_t)}^{2}\bigr]
\le
\frac{4\bigl(f(w_0)-\EE[f(w_R)]\bigr)}{\eta E R}
+ 2\zeta^{2}
+ \frac{4L\eta\sigma^{2}}{K}.
\tag{17}
\]

Since $f(w_R)\ge f^\star$, we may replace $\EE[f(w_R)]$ by $f^\star$ and obtain
\[
\frac{1}{R}\sum_{t=0}^{R-1}\EE\!\bigl[\norm{\nabla f(w_t)}^{2}\bigr]
\le
\frac{4\bigl(f(w_0)-f^\star\bigr)}{\eta E R}
+ 2\zeta^{2}
+ \frac{4L\eta\sigma^{2}}{K}.
\tag{18}
\]

\subsection*{Step 10: From average to a uniformly sampled iterate}
Select $\widetilde w$ uniformly from $\{w_t\}_{t=0}^{R-1}$. By definition,
\[
\EE\!\bigl[\norm{\nabla f(\widetilde w)}^{2}\bigr]
= \frac{1}{R}\sum_{t=0}^{R-1}\EE\!\bigl[\norm{\nabla f(w_t)}^{2}\bigr].
\]
Thus (18) directly yields the bound (T1) after multiplying numerator and denominator by $2$ to match the statement:
\[
\EE\!\bigl[\norm{\nabla f(\widetilde w)}^{2}\bigr]
\le
\frac{2\bigl(f(w_0)-f^\star\bigr)}{\eta RE}
+ L\eta\bigl(\sigma^{2}+E\zeta^{2}\bigr).
\tag{19}
\]

\subsection*{Step 11: Choice of $\eta$}
Set $\eta = \frac{c}{\sqrt{R}}$ with $c\le \frac{1}{L\sqrt{E}}$ (to respect (A4)).  
Plugging into (19) gives
\[
\EE\!\bigl[\norm{\nabla f(\widetilde w)}^{2}\bigr]
= O\!\Bigl(\frac{1}{\sqrt{R}}\Bigr)
= O\!\Bigl(\frac{1}{\sqrt{T/E}}\Bigr),
\]
which completes the proof. \qedhere

\section*{Rate Derivation (With Constants)}
From (19) let
\[
C_1 = 2\bigl(f(w_0)-f^\star\bigr),\qquad
C_2 = L\bigl(\sigma^{2}+E\zeta^{2}\bigr).
\]
Then
\[
\EE\!\bigl[\norm{\nabla f(\widetilde w)}^{2}\bigr]
\le \frac{C_1}{\eta RE}+C_2\eta .
\]
Minimising the right‑hand side w.r.t. $\eta$ yields $\eta^\star=\sqrt{C_1/(C_2 RE)}$, which (subject to (A4)) gives
\[
\EE\!\bigl[\norm{\nabla f(\widetilde w)}^{2}\bigr]\le
2\sqrt{\frac{C_1 C_2}{RE}} = O\!\Bigl(\frac{1}{\sqrt{R}}\Bigr).
\]

\section*{Special Cases}
\begin{itemize}
\item \textbf{Homogeneous data ($\zeta=0$).} The drift term disappears and (T1) reduces to the classic SGD rate $O(1/\sqrt{RE})$.
\item \textbf{Single local epoch ($E=1$).} The bound matches that of standard mini‑batch SGD with $K$ parallel workers.
\item \textbf{Large $E$ with bounded drift.} If $\zeta^{2}=O(1/E)$, the term $L\eta E\zeta^{2}$ remains $O(L\eta)$, preserving the $1/\sqrt{R}$ rate.
\end{itemize}

\section*{Discussion}
The theorem shows that FedAvg with $E$ local SGD steps converges to a stationary point of the global non‑convex loss at the same $O(1/\sqrt{R})$ rate as centralized SGD, provided the stepsize is chosen according to (A4) and the client heterogeneity is controlled via $\zeta$. The explicit drift term $\Delta_t$ appears only as an additive variance‑like component, confirming the intuitive trade‑off: more local computation (larger $E$) reduces communication but amplifies the effect of data heterogeneity.

\end{document}